\chapter{Edge-Host Connectivity: PoE, Voice, APs and IoT}\label{ch:edgehost}
\bp{2.2}

\begin{outcomes}
By the end of this chapter you will be able to:
\begin{tight}
  \item \textbf{Configure} an access port correctly for each of the five device
        types the blueprint names.
  \item \textbf{Explain} how a data VLAN and a voice VLAN share one cable, and
        why that is not a trunk in the ordinary sense.
  \item \textbf{Configure} and verify Power over Ethernet, including budget and
        per-port limits.
  \item \textbf{Diagnose} the faults peculiar to edge ports: a phone that gets no
        power, an AP that sees only one VLAN, and a device in the wrong VLAN
        because it was plugged into the wrong socket.
\end{tight}
\end{outcomes}

\begin{prereq}
\chref{vlans} for VLANs and access ports. \chref{cli} for interface ranges.
\end{prereq}

\begin{keyterms}
Power over Ethernet $\bullet$ PoE, PoE+, UPoE $\bullet$ power budget $\bullet$
voice VLAN $\bullet$ data VLAN $\bullet$ standalone AP $\bullet$ controller-based
AP $\bullet$ LWAPP/CAPWAP $\bullet$ IoT appliance $\bullet$ errdisable
\end{keyterms}

\begin{examnote}[Read topic 2.2 carefully --- it lists the devices]
Topic \tp{2.2} is ``configure Layer 2 switch port attributes for edge-host
connectivity (VLAN, PoE, port channel, and LACP)'' and then names five
categories: desktops, printers and IoT appliances; wireless access points,
standalone and controller-based; VoIP phones; virtualised hosts; and network
appliances. The exam expects a different port configuration for several of them,
so learn them as five cases rather than one.
\end{examnote}

\section{Power over Ethernet}

\begin{definitionbox}[PoE]
\term{Power over Ethernet} delivers electrical power over the same twisted pair
that carries data, so a phone, camera or access point needs one cable rather than
two. The switch negotiates with the device and supplies only what it asks for.
\end{definitionbox}

\begin{longtable}{@{}L{3.0cm}C{2.4cm}C{2.4cm}L{7.0cm}@{}}
\toprule
\rowcolor{primary!15}
\textbf{Standard} & \textbf{Per port} & \textbf{Name} & \textbf{Typically powers} \\
\midrule
\endhead
802.3af & 15.4 W & PoE & IP phone, basic AP, door sensor \\
\rowcolor{lightbg}
802.3at & 30 W & PoE+ & Modern AP, PTZ camera \\
802.3bt & 60--90 W & UPoE, PoE++ & Multi-radio AP, display, thin client \\
\bottomrule
\end{longtable}

\begin{alertbox}[The budget is the constraint, not the port]
A switch has a total power budget shared across all ports --- often far less than
the number of ports times the maximum per port. A 24-port PoE+ switch with a
370~W budget cannot deliver 30~W to all 24. When the budget is exhausted the
switch stops powering \emph{new} devices, so the symptom is ``the phone I plugged
in this morning works, the one I plugged in this afternoon does not'', and
nothing in the port configuration looks wrong.

Check \cmd{show power inline} before you check anything else.
\end{alertbox}

\begin{config}{PoE on an access port}
interface GigabitEthernet0/10
 description AP in the east corridor
 switchport mode access
 switchport access vlan 10
 ! auto is the default: negotiate and give the device what it asks for.
 power inline auto
 ! Cap a port that must never take more than a phone's worth.
 ! power inline auto max 15400
 no shutdown
!
! Never power this port -- a data-only port that must not energise.
interface GigabitEthernet0/24
 power inline never
\end{config}

\begin{verify}{SW1}{show power inline}
Available:370.0(w)  Used:71.5(w)  Remaining:298.5(w)

Interface Admin  Oper       Power   Device              Class Max
                            (Watts)
--------- ------ ---------- ------- ------------------- ----- ----
Gi0/1     auto   on         6.3     IP Phone 7960       2     15.4
Gi0/2     auto   on         6.3     IP Phone 7960       2     15.4
Gi0/10    auto   on         28.6    AIR-AP2802I         4     30.0
Gi0/11    auto   off        0.0     n/a                 n/a   30.0
Gi0/24    never  off        0.0     n/a                 n/a   0.0
\end{verify}

Read the top line first. \cmd{Remaining} is what decides whether the next device
you plug in will power up at all.

\section{Voice VLANs: two VLANs, one cable}

An IP phone has a small switch inside it: the cable from the wall goes into the
phone, and the PC plugs into the back of the phone. Both need network access, and
they should not be in the same VLAN --- voice needs its own subnet for quality
and for security.

\ccnafig{%
\begin{tikzpicture}[font=\scriptsize]
  \tikzset{
    dv/.style={rounded corners=3pt, draw=#1, fill=#1!8, text=#1!45!black,
               line width=0.9pt, minimum width=1.5cm, minimum height=1.0cm,
               align=center, font=\tiny\bfseries},
    w/.style={line width=0.9pt, neutral!65}}

  \node[dv=dom2, minimum width=1.7cm] (sw) at (0,0) {\faIcon{sitemap}\\SW1};
  \node[dv=dom5, minimum width=1.9cm] (ph) at (5.4,0) {\faIcon{phone-alt}\\IP phone};
  \node[dv=neutral] (pc) at (10.2,0) {\faIcon{desktop}\\PC};

  \draw[w, line width=1.6pt, dom2!55] (sw) -- (ph);
  \draw[w] (ph) -- (pc);

  \node[font=\tiny\bfseries, text=dom2!45!black, align=center] at (2.7,0.75)
      {\emph{one} cable from the wall};

  % The two flows. Labels sit ABOVE their arrows, not on top of them --- a
  % label with a white background wide enough to cover a short arrow erases it.
  \draw[-{Stealth[length=2mm]}, dom5!75, line width=0.9pt]
      (0.95,-1.15) -- (4.45,-1.15);
  \node[font=\tiny\bfseries, text=dom5!45!black, anchor=south] at (2.7,-1.05)
      {voice --- \textbf{tagged} VLAN 150};

  \draw[-{Stealth[length=2mm]}, neutral!70, line width=0.9pt]
      (0.95,-2.0) -- (9.45,-2.0);
  \node[font=\tiny\bfseries, text=neutral!55!black, anchor=south] at (5.2,-1.9)
      {data --- \textbf{untagged} VLAN 10, straight through the phone};

  \node[font=\tiny, text=dom2!45!black, align=center, text width=5.6cm,
        anchor=north] at (2.7,-2.35)
      {the switch tells the phone which VLAN to tag,\\over CDP or LLDP-MED};

  \node[draw=dom4!55, fill=dom4!5, rounded corners=3pt, line width=0.8pt,
        font=\tiny, align=center, text width=3.5cm, inner sep=4pt,
        text=dom4!40!black, anchor=north] at (10.4,-0.8)
      {\textbf{still an access port}\\--- not a trunk};
\end{tikzpicture}}%
{One cable, two VLANs, and no trunk. The phone tags its own traffic and passes
the PC's through untouched.}{voiceport}

\begin{conceptbox}[How one access port carries two VLANs]
The port is configured with a \term{data VLAN} and a \term{voice VLAN}. The
switch tells the phone, over CDP or LLDP-MED, which VLAN ID to tag its own
traffic with. The phone tags voice frames and passes the PC's frames through
untagged.

So the switch sees tagged voice traffic and untagged data traffic on one access
port. Strictly this is a trunk carrying two VLANs, but it is configured as an
access port and Cisco calls it a \term{multi-VLAN access port}. Do not configure
\cmd{switchport mode trunk} here --- the voice VLAN command does what is needed,
and a real trunk would change how the PC's untagged traffic is handled.
\end{conceptbox}

\begin{config}{a desk port for a phone with a PC behind it}
interface GigabitEthernet0/1
 description Desk 1 - phone with PC daisy-chained
 switchport mode access
 switchport access vlan 10
 switchport voice vlan 30
 ! The PC should forward immediately rather than waiting for
 ! spanning tree -- see chapter on Rapid PVST+.
 spanning-tree portfast
 power inline auto
 no shutdown
\end{config}

\begin{verify}{SW1}{show interfaces GigabitEthernet0/1 switchport}
Name: Gi0/1
Switchport: Enabled
Administrative Mode: static access
Operational Mode: static access
Access Mode VLAN: 10 (Staff)
Voice VLAN: 30 (Voice)
\end{verify}

\section{Access points, standalone and controller-based}

The blueprint distinguishes the two, and the distinction changes the port
configuration completely.

\begin{longtable}{@{}L{3.8cm}L{5.4cm}L{6.0cm}@{}}
\toprule
\rowcolor{primary!15}
\textbf{} & \textbf{Standalone (autonomous)} & \textbf{Controller-based (lightweight)} \\
\midrule
\endhead
Where the config lives & On the AP itself & On the wireless LAN controller \\
\rowcolor{lightbg}
What crosses the wire & Each SSID's traffic in its own VLAN & All client traffic
tunnelled to the controller in CAPWAP \\
Switch port & \textbf{Trunk}, carrying every SSID's VLAN & \textbf{Access port},
in the AP management VLAN \\
\rowcolor{lightbg}
Scales to & A handful of APs & Hundreds \\
\bottomrule
\end{longtable}

\begin{alertbox}[The most common AP misconfiguration]
Configuring a lightweight AP's port as a trunk, or a standalone AP's port as an
access port. Both leave the AP reachable --- it will get an address and appear
healthy --- while clients on some or all SSIDs get nothing. If wireless users on
one SSID work and another SSID does not, look at the switch port before you touch
the wireless configuration.
\end{alertbox}

\begin{config}{both kinds of access point}
! Lightweight AP: one access port in the AP management VLAN.
! Client traffic is tunnelled to the controller, so no client
! VLANs need to exist on this switch at all.
interface GigabitEthernet0/10
 description Lightweight AP - CAPWAP to WLC
 switchport mode access
 switchport access vlan 50
 power inline auto
 spanning-tree portfast
!
! Standalone AP: a trunk, because the AP bridges each SSID
! into a different VLAN and tags them itself.
interface GigabitEthernet0/11
 description Standalone AP - SSIDs in VLANs 10 and 20
 switchport trunk encapsulation dot1q
 switchport mode trunk
 switchport trunk native vlan 50
 switchport trunk allowed vlan 10,20,50
 power inline auto
\end{config}

\section{Printers, IoT appliances and virtualised hosts}

\begin{tight}
  \item \textbf{Printers and IoT appliances} --- ordinary access ports. Put them
        in their own VLAN: they are rarely patched, often speak plaintext
        protocols, and are a favourite foothold. Add port security
        (\chref{l2security}) because they never move.
  \item \textbf{Virtualised hosts} --- a server running virtual machines usually
        needs a \emph{trunk}, because each VM may sit in a different VLAN and the
        hypervisor's virtual switch tags accordingly. This is the one ``end
        device'' that is not an access port.
  \item \textbf{Network appliances} --- firewalls, load balancers and the like.
        Read the vendor's requirement; most want a trunk, some want several
        access ports.
\end{tight}

\begin{pitfall}
\begin{tight}
  \item \textbf{PortFast on a port that leads to a switch.} PortFast skips the
        spanning tree listening and learning states. On an edge port that is
        correct; on a port leading to another switch it can create a loop. Use
        BPDU guard alongside it (\chref{stp}).
  \item \textbf{Assuming the phone's PC port is in the voice VLAN.} It is not ---
        untagged traffic from the PC lands in the data VLAN, which is the whole
        point.
  \item \textbf{Ignoring the power class.} A class 4 device on a switch with only
        802.3af ports will not power up, however correct the configuration looks.
  \item \textbf{Trunking to a lightweight AP.} See the warning above.
\end{tight}
\end{pitfall}

\begin{breakfix}[Half the wireless users work; the other half get no address.]
Users on the \cmd{Staff} SSID are fine. Users on \cmd{Guest} associate to the
wireless network but never receive an IP address. The AP is reachable and its
configuration has not changed.

\begin{verify}{SW1}{show interfaces GigabitEthernet0/11 trunk}
Port        Mode         Encapsulation  Status        Native vlan
Gi0/11      on           802.1q         trunking      50

Port        Vlans allowed on trunk
Gi0/11      10,50
\end{verify}

\textbf{Diagnosis.} This is a standalone AP, correctly trunked, but the allowed
VLAN list carries 10 and 50 only. VLAN 20 --- the Guest SSID's VLAN --- is not
permitted on the trunk, so Guest clients associate at the radio (which the AP
handles alone) and then have no path to the DHCP server. Association succeeding
while addressing fails is the signature of a VLAN that does not reach the wire.

\textbf{Fix.} Add the missing VLAN to the trunk:
\cmd{switchport trunk allowed vlan add 20}. Note \cmd{add} --- without it you
replace the whole list and drop VLANs 10 and 50, taking down the users who were
working. Confirm with \cmd{show interfaces trunk}, then have a Guest client
renew.
\end{breakfix}

\begin{lab}{Wire an edge properly}{Packet Tracer}
\textbf{Topology.} One PoE switch; an IP phone with a PC behind it; a lightweight
AP; a printer.

\begin{tasks}
  \item Create VLANs 10 (Staff), 20 (Guest), 30 (Voice), 50 (AP-Mgmt) and 60
        (Printers).
  \item Configure the phone port with data VLAN 10, voice VLAN 30, PortFast and
        PoE. Verify both VLANs appear on the one port.
  \item Configure the AP port as an access port in VLAN 50 with PoE.
  \item Configure the printer port in VLAN 60 with PortFast.
  \item Run \cmd{show power inline} and record the budget, the used figure and
        the class of each powered device.
  \item Set \cmd{power inline never} on the printer port and observe what happens
        to a PoE device moved there.
  \item \textbf{Extension.} Change the AP to standalone: reconfigure its port as
        a trunk carrying VLANs 10, 20 and 50, and explain in one sentence why the
        two AP types need opposite port configurations.
\end{tasks}
\end{lab}

\begin{checkpoint}
\begin{tightnum}
  \item An IP phone powers up and works; a PC plugged into the back of it gets an
        address in VLAN 10. Which VLAN carries the phone's own traffic, and how
        did the phone learn to use it?
  \item A new AP will not power on. \cmd{show power inline} shows the port
        \cmd{auto/off} and \cmd{Remaining:4.2(w)}. What is wrong?
  \item Why does a lightweight AP need only one VLAN on its switch port when it
        serves four SSIDs?
\end{tightnum}
\end{checkpoint}

\begin{chaptersummary}
\begin{tight}
  \item PoE classes: 15.4~W af, 30~W at, up to 90~W bt. The shared \emph{budget}
        is what actually runs out, so read \cmd{show power inline} first.
  \item A voice VLAN plus a data VLAN on one access port lets a phone and a PC
        share a cable: the phone tags, the PC does not.
  \item Lightweight APs take an access port in the management VLAN; standalone
        APs take a trunk carrying every SSID's VLAN. Getting this backwards
        breaks some SSIDs and not others.
  \item Printers and IoT devices belong in their own VLAN with port security.
        Virtualised hosts usually need a trunk.
  \item Use \cmd{switchport trunk allowed vlan add}, never bare \cmd{allowed
        vlan}, on a live trunk.
\end{tight}
\end{chaptersummary}

\begin{reviewq}
  \item \textbf{(MCQ)} Which port configuration suits a controller-based access
        point? (a)~trunk with all client VLANs (b)~access port in the AP
        management VLAN (c)~routed port (d)~access port in VLAN 1
  \item \textbf{(MCQ)} Which command stops a port ever supplying power?
        (a)~\cmd{power inline auto} (b)~\cmd{power inline never}
        (c)~\cmd{no power inline} (d)~\cmd{shutdown}
  \item \textbf{(Short)} Explain why a port with a voice VLAN is not configured
        as a trunk, even though it carries two VLANs.
  \item \textbf{(Short)} A device on a PoE port shows \cmd{Oper: off} while other
        ports are powering normally. Give two causes and the command that
        distinguishes them.
  \item \textbf{(Applied)} You must connect twelve IP phones with PCs behind
        them, four PoE+ access points and six printers to one 24-port switch with
        a 370~W budget. Estimate whether the budget suffices, state your
        assumptions, and write the configuration for one port of each type.
\end{reviewq}
